% !TeX root = ../教材题库_全部条目_转写工作稿.tex
% Source: XJ-B1-C04-EX-P116-Y0405
% 第4章 / 4.3 对数函数 / 教材 p.116 / PDF 第123页
% OCR未校对：公式与文字均可能有误；未完成数学TeX转写；禁止直接发给学生
% 原题及相邻原解截图（核对用）：../题目原图/XJ-B1-C04-EX-P116-Y0405-01.webp
% 本文件只在转写完成后作为题干使用；OCR 条目尚不是可靠题干。
\begin{bankitem}{4.3 例题 例题2}{第4章 / 4.3 对数函数 / 教材 p.116 / PDF 第123页}{XJ-B1-C04-EX-P116-Y0405}{OCR未校对：公式与文字均可能有误；未完成数学TeX转写；禁止直接发给学生}
\begin{ocrdraft}
[教材p.116 / PDF第123页]
\par
00求下列各式的值：
\par
logz0；
\par
（2）logo.61；
\par
23一2
\par
（4）2（3十2）
\par
（3）
\par
解
\par
（1）logz一一10g22一
\par
（2）
\par
logo61一log“0．6=0；
\par
23
\par
3
\par
（3）
\par
23一2一23\ensuremath{\cdot}2一
\par
22
\par
4
\par
（4）
\par
2嵊2（3十2）=25=5
\par
本节“思考”中，0．568J=0．01，则的值可表示为x
\par
一logo．5680．01
\par
，利用计
\par
算器或计算机等，可求得x一logo．5680．01\&142．
\end{ocrdraft}
\end{bankitem}
